\section{Positive boxes and coefficient regularity}\label{sec:positive-boxes}
Let $C_{\rm box}(\rho_B)$ be the supremal gap ratio for positive monomial
factorizations on strictly positive boxes with every coordinate aspect
ratio $u_i/\ell_i\le\rho_B$, over arbitrary dimensions and evaluation
points with positive hull gap. Unlike $C_\delta$, this restricts the box,
not the normalized evaluation point. Fixed coordinates may be retained for
the dimension-independent bound; remove them when counting the nonfixed
coordinates in the finite-dimensional refinement.

\subsection{A coefficient inequality with one ambient orientation law}
Fix an ambient dimension $N\ge2$. Choose a uniform subset $J\subseteq[N]$
of size $\lfloor N/2\rfloor$ and an independent uniform $U\in[0,1]$.
Coordinate $i$ of mean $p_i$ succeeds on $[0,p_i]$ if $i\in J$ and on
$[1-p_i,1]$ otherwise. Call this balanced orientation law $O$. It has
all prescribed means, though for odd $N$ its individual orientation
choices are not fair. The probability that a fixed pair has opposite
orientations is
\[
 q_N=\frac{2\lfloor N/2\rfloor\lceil N/2\rceil}{N(N-1)},\qquad
 \beta_N=q_N^{-1}=
 \begin{cases}2-2/N,&N\text{ even},\\2-2/(N+1),&N\text{ odd}.
 \end{cases}
\]
All local supports and smaller-dimensional induction steps below use
restrictions of this one law. Deleting a deterministic coordinate
neither conditions nor resamples the orientation subset.

For a support $S$ of size $d\le N$, let $K=\sum_{i\in S}X_i$.
Write $C_j,P_j,O_j$ for $\E\binom Kj$ under common-threshold,
independent, and ambient orientation rounding. With
$s=\sum_{i\in S}p_i$, $k=\lfloor s\rfloor$, and $\theta=s-k$, put
\[
 V_j=(1-\theta)\binom kj+\theta\binom{k+1}j.
\]
Degree-zero moments are one; negative orders and orders above support
size are zero. By Lemma~\ref{lem:cardinality-envelopes}, $C_j$ and $V_j$
are the true upper and lower envelope values of the elementary
symmetric polynomial of degree $j$.

\begin{lemma}[Finite spreading inequality]\label{lem:box-coefficient}
For every support $S\subseteq[N]$, every mean vector, and every integer
$j\ge0$,
\begin{equation}\label{eq:box-coefficient}
 P_j-V_j\le\beta_N(C_j-O_j)+(C_{j-1}-P_{j-1}).
\end{equation}
The same statement holds with $\beta_N$ replaced by $2$ and independently
fair endpoint orientations in every coordinate.
\end{lemma}
\begin{proof}
Write $\beta$ for the indicated constant and set
$F_{S,j}=\beta C_j+V_j-P_j-\beta O_j+C_{j-1}-P_{j-1}$.
Orders zero and one vanish; order $d+1$ is $C_d-P_d\ge0$, and higher
orders vanish. For order two let
$L_2=\sum_{a<b}\max(0,p_a+p_b-1)$. Equal orientations give a pair's
upper Fr\'echet value; opposite orientations give its lower value.
Their opposite probability is $1/\beta$, including under restriction.
Thus $O_2=(1-1/\beta)C_2+L_2/\beta$, and
\[
 F_{S,2}=(C_2-P_2)+(V_2-L_2)\ge0.
\]
The first inequality follows from upper optimality, the second because
every law, including an adjacent-cardinality law, respects every pair
lower bound.

We induct on support size. A mean-zero coordinate is deterministic and
its deletion gives $F_{S,j}=F_{S\setminus\{i\},j}$.
A mean-one coordinate and Pascal's identity give
$F_{S,j}=F_{S\setminus\{i\},j}+F_{S\setminus\{i\},j-1}$.
The adjacent-cardinality values obey the same identity after shifting
the sum by one. These are restrictions of the ambient law, not newly
sampled balanced laws, so its correlated coins cause no exception.
All boundary points follow by induction; dimension one is immediate.

For an interior vector and $3\le j\le d$, select distinct coordinates
of global minimum $x$ and global maximum $y$. Move them to $x-h,y+h$
with $0\le h\le\min(x,1-y)$. The sum stays fixed, so $\Delta V_j=0$.
In increasing order of means,
\[
 C_j=\sum_{i=1}^d p_{(i)}\binom{d-i}{j-1}.
\]
Putting the selected minimum first and maximum last also handles ties.
Hence, with $B=\binom{d-1}{j-1}$ and $D=\binom{d-1}{j-2}$,
$\Delta C_j=-Bh$ and $\Delta C_{j-1}=-Dh$; the last coordinate has
zero coefficient since $j\ge3$. If $E_r$ is the elementary symmetric
polynomial in the other $d-2$ means, the unchanged sum and product
change $\Delta(xy)=-h(y-x+h)$ give exactly
\[
 \Delta(-P_j-P_{j-1})
 =h(y-x+h)(E_{j-2}+E_{j-3})\le Dh.
\]
Here $0\le y-x+h\le1$ and $E_r\le\binom{d-2}{r}$.
Finally couple both orientation vectors using the same $U$ and coins.
Increasing a single mean by $h$ changes only that bit, monotonically,
on an event of probability $h$. The increase in $\binom Kj$ is then
$\binom{K_{\rm rest}}{j-1}\le B$. Decreasing the minimum first and
increasing the maximum second therefore gives $\Delta O_j\ge-Bh$,
regardless of dependence among orientation coins. Combining the finite
changes yields $\Delta F_{S,j}\le-\beta Bh-Dh+Dh+\beta Bh=0$.
Take $h=\min(x,1-y)$ to reach a boundary point, whose coefficient is
nonnegative by induction. The original coefficient is at least that
boundary value. No differentiability assumption is used.
\end{proof}

The choice of global extrema in this proof matters. For independently
fair orientations, the stronger Schur-concavity assertion is false:
\begin{equation}\label{eq:schur-counterexample}
 \begin{split}
 F_{[4],3}(2/5,7/10,24/25,97/100)&=6201/12500,\\
 F_{[4],3}(2/5,7/10,959/1000,971/1000)&=4961031/10000000.
 \end{split}
\end{equation}
Spreading the last two coordinates by $1/1000$ increases the coefficient
by $231/10000000$. These rational values follow by averaging the sixteen
orientation patterns, integrating between their interval endpoints,
and evaluating the displayed symmetric moments. They refute arbitrary
pair spreading, without affecting the global-minimum/global-maximum
comparison.

\begin{theorem}[Positive-box upper bounds]\label{thm:positive-box-upper}
For $N\ge2$ nonfixed coordinates and $\rho_B>1$, every positive
multilinear polynomial on a strictly positive box of aspect ratios at
most $\rho_B$ satisfies
\[
 \tgap\le(\rho_B+\beta_N)\hgap\le(\rho_B+2)\hgap.
\]
The common mixture of independence and balanced orientation with weights
$\rho_B/(\rho_B+\beta_N)$ and $\beta_N/(\rho_B+\beta_N)$ captures this
fraction of each original monomial's exact gap simultaneously. For zero
or one nonfixed coordinate all gaps vanish.
\end{theorem}
\begin{proof}
First consider a support on $[1,\rho_B]^N$ and put $t=\rho_B-1$.
Its binary physical value is $(1+t)^K$, so its upper, independent,
orientation, and lower values are $C(t)=\sum_jC_jt^j$, $P(t)$,
$O(t)$, and $V(t)=\sum_jV_jt^j$. The last identity follows from a
single adjacent-cardinality law attaining every order together.
The coefficient of $t^j$ in
$(\beta_N+t)C+V-(1+t)P-\beta_NO$ is $F_{S,j}\ge0$.
Consequently the true local gap satisfies
\begin{equation}\label{eq:physical-local}
 C-V\le\rho_B(C-P)+\beta_N(C-O).
\end{equation}
The same mixture works on all supports in the ambient $N$ coordinates.
Summing with nonnegative coefficients and using common upper attainment
proves the common-aspect conclusion. Fair independent orientation coins
and $\beta=2$ prove $\rho_B+2$ directly in every dimension.

For unequal endpoints, including fixed coordinates, let $0<\ell_i\le u_i$ and put
\[
 s_i=\frac{u_i/\ell_i-1}{\rho_B-1}\in[0,1],\qquad
 z_i=\ell_i[(1-s_i)+s_i y_i],\quad y_i\in[1,\rho_B].
\]
This is an affine surjection onto the original box. A fixed coordinate
has $s_i=0$ and disappears from the expanded function. Pushing laws
forward preserves objective expectations; conversely, lift a law through
the inverses on nonfixed coordinates and assign the unused coordinates
the chosen preimage means. Thus the full hull gap is preserved at every
preimage of the evaluation point.

Each original monomial becomes a nonnegative sum of physical
$y$-monomials. Its concave envelope equals the sum of the expanded
concave envelopes by common-threshold attainment. Its convex envelope
is at least the sum of the expanded convex envelopes, so its true
factor gap is at most the sum of their gaps.
Apply the common-aspect bound to the expansion, always using restrictions
of the same ambient orientation law. To obtain $\beta_N$ with $N$ counting
only nonfixed coordinates, first discard fixed coordinates and select
the ambient law in those $N$ coordinates. The equality of concave values
also makes the original factor's deficiency under either law the sum of
expanded deficiencies, proving the claimed per-original-factor guarantee.
The certificate can be sampled without forming the expansion: the
endpoint laws pull back to the original endpoints. Box inclusion alone
would not prove any of these comparisons.
\end{proof}

The repository's
\href{../formal/topics/18-positive-box/COVERAGE.md}{topic-18 coverage map}
links the coefficient inequality, original-box transfer and balanced
refinement to their Lean declarations. Its
\href{../formal/topics/18-positive-box/REVIEW.md}{review} records the
resolved findings, and its
\href{../formal/topics/18-positive-box/VERIFICATION.md}{verification record}
states which builds and kernel replays were completed. These results
establish the stated upper bounds without claiming minimality of $\beta_N$.

\begin{theorem}[Coefficient-regularity cardinality bound]
\label{thm:coefficient-regularity}
In ambient dimension $N\ge2$, suppose each entire local factor is the
multiaffine interpolant of
\[
 \psi_e(k)=\sum_{j=0}^{d_e}a_{e,j}\binom kj,\qquad a_{e,j}\ge0,
 \qquad a_{e,j+1}\le L a_{e,j}\quad(j\ge2),
\]
where $L\ge0$ is common to all factors and coefficients beyond the
degree are zero. Then
\[
 \tgap\le(L+1+\beta_N)\hgap.
\]
Separate affine terms may be added freely. Fair independent endpoint
orientations give the dimension-free bound $L+3$. These are bounds
for the gaps of the entire cardinality factors, not the gaps obtained
by separately relaxing their expanded monomials.
\end{theorem}
\begin{proof}
The nonnegative binomial coefficients make each sequence discrete convex:
its second difference is $\sum_{j\ge2}a_{e,j}\binom{k}{j-2}\ge0$.
Lemma~\ref{lem:cardinality-envelopes} supplies a common upper extremizer,
and within each factor one adjacent-cardinality law attains every
$V_j$ together. Denote that factor's four values by $C_\psi,V_\psi,
P_\psi,O_\psi$. Multiply \eqref{eq:box-coefficient} by its $a_j$ and
sum. Since $C_j-P_j\ge0$ and the order-zero and order-one deficiencies
vanish, reindexing gives
\[
 \begin{split}
 P_\psi-V_\psi
 &\le\beta_N(C_\psi-O_\psi)
          +\sum_{j\ge2}a_{j+1}(C_j-P_j)\\
 &\le\beta_N(C_\psi-O_\psi)+L(C_\psi-P_\psi).
 \end{split}
\]
Add $C_\psi-P_\psi$ to obtain
\[
 C_\psi-V_\psi\le(L+1)(C_\psi-P_\psi)+\beta_N(C_\psi-O_\psi).
\]
One global mixture with weights proportional to $L+1$ and $\beta_N$
therefore captures the required fraction for all factors. Summing
and using common upper attainment proves the theorem. This also
handles $L=0$, deterministic means, and affine factors. With at most
one nonfixed coordinate the interpolants are affine and all gaps vanish.
The product sequence $\rho_B^k$ has $a_j=(\rho_B-1)^j$, so
$L=\rho_B-1$ recovers the common-aspect product inequality.
\end{proof}
This refinement is a direct consequence of the coefficient mechanism;
no optimality or independent priority is asserted. Distinct earlier
low/high and asymmetric proofs, and the precise limitation of fixed
independence/orientation mixtures, appear in
Appendix~\ref{app:positive-box-predecessors}.

\subsection{A lower family on each fixed positive box}
\begin{theorem}[Aspect-ratio lower bound]\label{thm:positive-box-lower}
For every $\rho_B>1$,
\[
 \max\{2,\rho_B\}\le C_{\rm box}(\rho_B)\le\rho_B+2.
\]
Thus $C_{\rm box}(\rho_B)=\rho_B+O(1)$ as $\rho_B\to\infty$.
The exact finite-aspect-ratio constant is not determined.
\end{theorem}
\begin{proof}
The lower bound two already holds at width two by
Theorem~\ref{thm:width-two-gap}. For the lower bound $\rho_B$, fix
$\epsilon=1/\rho_B\in(0,1)$ and put $\alpha=1-\epsilon$.
Use the radix supports with $L\ge2$, $b=L^2$, $m=b^L$, now on
$[\epsilon,1]^{m+L}$ with normalized means
$\E A_j=b^{-j}$ and $\E Z_i=1-m^{-1}$. A level-$j$ term has
$k=m/b^j$ leaves and total failure mean $(1-b^{-j})+k/m=1$.
Its vertex value is $\epsilon^{R_e}$, so its exact lower value is
$\epsilon$, attained by exactly one failure with the prescribed
categorical probabilities. Common-threshold rounding gives its upper
value
\[
 C_{j,Q}=\epsilon+\alpha b^{-j}
                      -\frac\epsilon m(1-\epsilon^k).
\]
Indeed its successive threshold states have product values
$1,\epsilon,\epsilon^{k+1}$ with interval lengths
$b^{-j},1-m^{-1}-b^{-j},m^{-1}$. Therefore
\begin{equation}\label{eq:positive-radix-term}
 \tgap_L=\alpha L-D_L,\qquad
 D_L=\epsilon\sum_{t=0}^{L-1}\frac{1-\epsilon^{b^t}}{b^t},
 \qquad0\le D_L\le\frac{\epsilon b}{b-1}.
\end{equation}
For any common endpoint law let $R_Q$ count failed leaves in block $Q$
and $R$ count them in total. Since $\E R=1$, subtracting expected
physical products from their common upper value gives exactly
\begin{equation}\label{eq:softened-coverage}
 \hgap_L=\max\E\left[
 \alpha\sum_jA_j\sum_{Q\in\mathcal P_j}(1-\epsilon^{R_Q})
 +\epsilon\sum_j\sum_{Q\in\mathcal P_j}(1-\epsilon^{R_Q})
 \right]-D_L.
\end{equation}
For its first term use $1-\epsilon^{R_Q}\le\ind\{R_Q>0\}$ and the
unit-cube radix hull upper bound. For its second use
$1-\epsilon^r\le\alpha r$ and $\sum_{Q\in\mathcal P_j}R_Q=R$.
This proves
\[
 \hgap_L\le\epsilon\alpha L+
                 \alpha[1+(L-1)/b]-D_L.
\]
Exactly one uniformly failed leaf, with independent anchors of the
required means, instead gives
\[
 \hgap_L\ge\epsilon\alpha L+
                    \alpha^2\sum_{j=1}^Lb^{-j}-D_L.
\]
The same geometric inequality gives $D_L\le\epsilon\alpha L$, so this
last bound is strictly positive for every finite $L$. The bounded
correction in \eqref{eq:positive-radix-term} now implies
$\tgap_L/L\to\alpha$, $\hgap_L/L\to\epsilon\alpha$, and their ratio
tends to $\rho_B$. Scale coordinates by $\rho_B$ to get positive
coefficients on $[1,\rho_B]^{m+L}$; both gaps are preserved. This is
a limit in dimension for each fixed aspect ratio, requiring no exchange
of those limits. The upper bound is Theorem~\ref{thm:positive-box-upper}.
\end{proof}

\subsection{Unequal aspect ratios at frequency two}
\begin{example}[A rational bipartite ratio $7/6$]\label{ex:unequal-seven-six}
For $f(x,y,z)=xy+xyz$ on $[1,2]\times[1,2]\times[1,3]$, evaluate at
$(5/4,5/4,5/2)$. Put $x=1+a,y=1+b,z=1+2c$, so the normalized mean is
$(1/4,1/4,3/4)$. The following valid affine bounds give exact values:
\[
 xy\ge1+a+b,\quad xyz\ge2a+2b+3c,\quad
 f\ge4a+4b+4c,\quad f\le2+8a+4b+2c.
\]
Validity follows by substitution at all eight binary vertices and
multiaffine interpolation. The respective attaining distributions are
\[
\begin{array}{c|l}
 xy\text{ lower}&000,001,011,101\text{ each of mass }1/4\\
 xyz\text{ lower}&001\text{ of mass }3/4;\ 110\text{ of mass }1/4\\
 f\text{ lower}&001\text{ of mass }1/2;\ 011,100\text{ each of mass }1/4\\
 f\text{ upper}&000\text{ of mass }1/4;\ 001\text{ of mass }1/2;
                         \ 111\text{ of mass }1/4.
\end{array}
\]
Each row has the required means. The separate upper bounds
$xy\le1+2a+b$ and $xyz\le1+6a+3b+2c$ are attained by the last law.
Thus $\tgap=13/2-19/4=7/4$, while $\hgap=13/2-5=3/2$.
The dual multigraph is bipartite: its two factor vertices share two
parallel edges, with a dummy leaf for $z$. This disproves unequal-box
bipartite scalar exactness, while leaving the common-aspect theorem intact.
\end{example}

\begin{proposition}[Bipartite positive-box boundary]\label{prop:unequal-bipartite}
For arbitrary nonnegative boxes with a bipartite dual graph, positive
monomial gaps satisfy $\tgap\le2\hgap$. Their worst ratio over strictly
positive boxes is at least $3/2$, already with two factors and three
variables. Its exact value between $3/2$ and $2$ remains open here,
as does a universal $3/2$ bound for unequal-box frequency-two instances.
\end{proposition}
\begin{proof}
More generally a proper coloring of the term conflict graph into $q$
classes gives $\tgap\le q\hgap$: terms in one class have disjoint
scopes, so join their minimizing laws independently and complete unused
means. They attain their entire class gap, and other factors have
nonnegative expected deficiency from their upper envelopes. Since those
upper values are commonly attained, $\hgap$ is at least each class's
total gap. Sum these inequalities. This also proves a bound two for
any two positive terms regardless of their boxes or frequency.

For the lower family let $0<\epsilon\le2$ and use
$f_\epsilon=(2/\epsilon)xy+xyz$, with
$x,y\in[1,1+\epsilon]$, $z\in[1,3]$, at normalized means
$(a,b,c)=(1/4,1/4,3/4)$. After removing affine parts the two factors are
\[
 A=2\epsilon ab,\qquad
 B=\epsilon^2ab+2\epsilon ac+2\epsilon bc+2\epsilon^2abc,
 \qquad C=A+B.
\]
Their lower bounds at the point are attained at the first three laws
of Example~\ref{ex:unequal-seven-six}, respectively. The affine
certificates are $A\ge0$, $B\ge\epsilon^2(a+b+c-1)$, and
$C\ge2\epsilon(a+b+c-1)$. In binary order
$000,001,010,011,100,101,110,111$, the last two residual lists are
\[
 \begin{split}
 &(\epsilon^2,0,0,\epsilon(2-\epsilon),0,
                  \epsilon(2-\epsilon),0,\epsilon(\epsilon+4)),\\
 &(2\epsilon,0,0,0,0,0,\epsilon^2,\epsilon(3\epsilon+2)),
 \end{split}
\]
which are nonnegative throughout the stated range. The local lower
values sum to $\epsilon^2/4$ and the full lower value is $\epsilon/2$.
The same common upper law as before gives $3\epsilon/2+3\epsilon^2/4$.
Hence
\[
 \tgap=\epsilon(3/2+\epsilon/2),\quad
 \hgap=\epsilon(1+3\epsilon/4),\quad
 \frac{\tgap}{\hgap}=\frac{2(\epsilon+3)}{3\epsilon+4}\longrightarrow3/2.
\]
To obtain unit coefficients divide the polynomial by $2/\epsilon$ and
put $z'=(\epsilon/2)z$. It becomes $xy+xyz'$ on
$[1,1+\epsilon]^2\times[\epsilon/2,3\epsilon/2]$, with the same ratio.
\end{proof}
For frequency-two products on unequal positive boxes, normalizing upper
endpoints to one gives a binary failure cost
$\prod_{i\in J\cap e}\alpha_i=\exp(-\sum_{i\in J\cap e}w_i)$,
where $\alpha_i=\ell_i/u_i$ and $w_i=-\log\alpha_i$. Each edge has
the \emph{same} weight $w_i$ at both incident factors. This is a weighted
convex degree potential, whereas the matching theorem above uses
unweighted cardinalities. Finite floating-point searches approaching
$3/2$ from below do not supply an upper bound. If the dual multigraph
is a forest, conditional gluing actually preserves all local minima
and maxima even for signed multiaffine factors on arbitrary boxes;
Appendix~\ref{app:structural-auxiliary} gives the elementary argument.
